\chapter{Overview to Particle Physics}

\section{Elementary particles and interactions: the energy-scale viewpoint}
The basic questions of particle physics are deceptively simple:
\begin{itemize}
  \item What are the elementary building blocks of Nature?
  \item What are the interactions among them?
\end{itemize}
A key lesson is that the answer depends on \emph{the energy scale} at which we probe Nature:
as we increase the typical energy (or momentum transfer) of our experiments, new degrees of
freedom become relevant and the effective description changes.

A pragmatic (and historically successful) way to organize our knowledge is to relate
\emph{energy scales} to the \emph{particles we can resolve}, the \emph{interactions we observe},
their \emph{ranges}, and their \emph{mediators}. At very low energies we describe matter in terms
of atoms and effective potentials; at higher energies we resolve nuclei, then nucleons,
and eventually quarks, leptons and gauge bosons.

\subsection{A qualitative ``energy ladder''}
\begin{center}
\small
\setlength{\tabcolsep}{8pt}
\renewcommand{\arraystretch}{1.15}
\begin{tabular}{@{}c c c c@{}}
\toprule
\textbf{Scale} & \textbf{Degrees of freedom} & \textbf{Interactions} & \textbf{Range} \\
\midrule
$\lesssim \mathrm{eV}$ &
atoms (effective) &
effective (atomic/molecular) &
long \\

$\sim \mathrm{keV}$ &
$e^-$, nuclei &
e.m. &
long \\

$\sim \mathrm{MeV}$ &
$e^\pm$, $p,n$, $\nu_e$ &
e.m.\ + nuclear + weak &
long / short / contact \\

$\sim \mathrm{GeV}$ &
$\mu^\pm$, $\nu_\mu$, hadrons &
e.m.\ + strong + weak &
long / short / contact \\

$\sim 10~\mathrm{GeV}$ &
$\tau^\pm$, $\nu_\tau$, $u,d,s,c,b$ &
e.m.\ + QCD + weak &
long / short / contact \\

$\sim \mathrm{TeV}$ &
$W^\pm$, $Z^0$, $t$, $H$ &
e.m.\ + QCD + EW &
long / short / short \\
\bottomrule
\end{tabular}
\end{center}



\section{Natural units and physical intuition}
At subatomic scales we inevitably need quantum mechanics, hence $\hbar$ appears in formulas.
At high velocities we need special relativity, hence $c$ appears as well. To avoid a
constant proliferation of $\hbar$ and $c$, we adopt \emph{natural units}
\begin{equation}
\hbar = c = 1,
\end{equation}
so that any dimensionful quantity can be expressed as a power of energy. In this convention,
masses, momenta and energies share the same units (eV, MeV, GeV, \dots).

To recover SI units, one reinserts appropriate powers of $\hbar$ and $c$ and uses
$1~\text{eV} \simeq 1.6\times 10^{-19}\,\text{J}$. Useful reference scales include
\begin{itemize}
  \item $E_{\text{Rydberg}}=13.6~\text{eV}$,
  \item $m_e\simeq 0.5~\text{MeV}$,
  \item $m_p\simeq 940~\text{MeV}$.
\end{itemize}
For distances and cross sections, it is common to use
\begin{itemize}
  \item the fermi (fm): $1~\text{fm}=10^{-15}\,\text{m}$, and $\hbar c\simeq 197~\text{MeV}\,\text{fm}$,
  \item the barn (b): $1~\text{b}=100~\text{fm}^2$.
\end{itemize}

\section{The Standard Model and what lies beyond}
The Standard Model (SM) describes elementary particles and their interactions up to energies
of order $\mathcal{O}(10~\text{TeV})$. It contains:
\begin{itemize}
  \item \textbf{Matter fields} (spin $J=\tfrac{1}{2}$): quarks and leptons. Charged leptons and quarks are massive; neutrino masses are very small.
  \item \textbf{Gauge bosons} (spin $J=1$): photon and gluons are massless; $W^\pm$ and $Z^0$ are massive.
  \item \textbf{Higgs boson} (spin $J=0$), massive.
\end{itemize}

Despite its success, there are clear indications that the SM is not the final story:
\begin{itemize}
  \item \textbf{Gravity:} extremely weak at subatomic distances but long-range, suggesting a massless mediator (graviton) in a quantum description.
  \item \textbf{Dark matter:} gravitational evidence at galactic/cosmological scales with no electromagnetic signals, suggesting at least one new particle.
\end{itemize}

\section{Energy frontiers in Nature}
The highest man-made energies are achieved in colliders. At present, the LHC reaches
$\sim 13~\text{TeV}$ in the center-of-mass energy. Future proposals include FCC(hh) at
$\sim 100~\text{TeV}$ and lepton/hadron options at intermediate energies.

Nature also provides \emph{cosmic rays} with energies far beyond colliders:
ultra-high energy cosmic rays satisfy $E>10^4~\text{TeV}$, with record events around
$E\sim 3\times 10^8~\text{TeV}$ (Utah, 1991).

\section{QFT as the language of particle physics}
Combining quantum mechanics with special relativity naturally leads to
\emph{quantum field theory} (QFT):
\begin{itemize}
  \item states can contain an \emph{arbitrary number of particles} (Fock space),
  \item elementary particles correspond to \emph{unitary irreducible representations} of the
  Poincar\'e group ISO$(3,1)$.
\end{itemize}
The irreducible unitary representations are characterized by
\begin{equation}
m^2 \in \mathbb{R}_+,\qquad \vec p\in \mathbb{R}^3,\qquad J\in \tfrac{1}{2}\mathbb{Z}_+,\qquad
J_3\in\{-J,-J+1,\dots,J\},
\end{equation}
with the special feature that for massless particles only the extreme spin projections appear.

\section{Strength of interactions and unification scales}
A convenient way to compare interactions is through dimensionless couplings (or effective
dimensionless combinations at a given scale). At $E\sim m_p\sim 1~\text{GeV}$ (low energies)
and $E\sim m_Z\sim 100~\text{GeV}$ (high energies), one finds qualitatively:
\begin{itemize}
  \item \textbf{Strong:} $\alpha_s(m_p)\sim 1$ but $\alpha_s(m_Z)\sim 0.1$ (running).
  \item \textbf{Electromagnetic:} $\alpha(m_p)\simeq 1/137$, $\alpha(m_Z)\simeq 1/129$.
  \item \textbf{Weak:} at low energies appears through $G_F m_p^2\sim 10^{-5}$, while at $m_Z$ one may define $\alpha_W(m_Z)=\dfrac{m_W^2 G_F}{4\pi}\sim \mathcal{O}(1)$.
  \item \textbf{Gravity:} $G m_p^2\sim 10^{-39}$, negligible in particle physics.
\end{itemize}

Two important conceptual scales follow from extrapolations:
\begin{itemize}
  \item \textbf{Electroweak unification:} around $E\sim m_Z$, electromagnetic and weak strengths become comparable.
  \item \textbf{Grand unification:} strong and electroweak strengths may meet around $E\sim 10^{16}~\text{GeV}$.
  \item \textbf{Planck scale:} gravity becomes comparable to other interactions at $E\sim G^{-1/2}\sim 10^{19}~\text{GeV}$.
\end{itemize}

\section{Conservation laws and discrete symmetries}
Conservation laws are deeply tied to symmetries:
\begin{itemize}
  \item \textbf{Space-time symmetries:} energy--momentum conservation; angular momentum conservation.
  \item \textbf{Internal symmetries:} electric charge; baryon number $B$; lepton number $L$.
\end{itemize}
In the SM, individual lepton flavor numbers (electron, muon, tau) are conserved at the Lagrangian level,
but neutrino oscillations show this is not an exact symmetry in Nature.
Quantum effects also violate $B$ and $L$ separately at an unobservably small level, leaving $B+L$
as the truly conserved combination within the SM description at accessible energies.

Strong and electromagnetic interactions conserve additional discrete symmetries:
parity (P), charge conjugation (C), time reversal (T), and (approximately) flavor.

\section{Hadrons: mesons, baryons and color}
Although quarks and gluons are the building blocks of the strong interaction in the SM,
in Nature we observe only \emph{hadrons}: bound states that are color singlets.
The simplest hadrons are:
\begin{itemize}
  \item \textbf{Mesons:} quark--antiquark ($q\bar q$),
  \item \textbf{Baryons:} three quarks ($qqq$).
\end{itemize}

\subsection{Isospin and light mesons/baryons}
Many light hadrons can be organized into representations of SU(2) isospin.
For example, pions form an isotriplet $(\pi^+,\pi^0,\pi^-)$ and nucleons form an isodoublet $(p,n)$.
A simple quark-model picture for light mesons treats them as $q\bar q$ with
\begin{equation}
J=\tfrac{1}{2}\otimes \tfrac{1}{2}=0\oplus 1,\qquad
I=\tfrac{1}{2}\otimes \tfrac{1}{2}=0\oplus 1,
\end{equation}
with $u$ and $d$ quarks carrying charges $Q_u=\tfrac{2}{3}e$ and $Q_d=-\tfrac{1}{3}e$.

Similarly, light baryons can be described as three-quark states with vanishing orbital
angular momentum:
\begin{equation}
\tfrac{1}{2}\otimes \tfrac{1}{2}\otimes \tfrac{1}{2}
= \tfrac{1}{2}\oplus \tfrac{1}{2}\oplus \tfrac{3}{2},
\end{equation}
leading to, for instance, the nucleon doublet and the $\Delta$ quartet.

\subsection{The Pauli principle and the need for color}
A puzzle arises because the spin-$3/2$ $\Delta$ states are totally symmetric under
exchange of quarks, apparently conflicting with the Pauli principle.
The resolution is the introduction of a new quantum number: \textbf{color}.
Quarks transform in the fundamental representation $\mathbf{3}$ of SU(3)$_\text{color}$, and
physical states must be color singlets.

Group-theory products show why mesons and baryons can be colorless, while isolated quarks
and gluons cannot appear as asymptotic states:
\begin{equation}
\mathbf{3}\otimes \mathbf{3}^\ast = \mathbf{1}\oplus \mathbf{8},\qquad
\mathbf{3}\otimes \mathbf{3} = \mathbf{3}^\ast \oplus \mathbf{6},\qquad
\mathbf{3}\otimes \mathbf{3}\otimes \mathbf{3} = \mathbf{1}\oplus \mathbf{8}\oplus \mathbf{8}\oplus \mathbf{10}.
\end{equation}
Hence $q\bar q$ and $qqq$ can contain singlets, but single quarks, gluons (octet), or $qq$ states
cannot be observed as free particles.

\subsection{Exotic hadrons}
QCD does not forbid hadrons beyond mesons and baryons: tetraquarks and pentaquarks are allowed.
Their unambiguous identification is easier in systems with heavy quarks, where the spectrum is
cleaner and thresholds are better controlled.

At low energies, the strong interaction is effectively short-range because the lightest
physical excitations exchanged between hadrons are pions, with $m_\pi \sim 140~\text{MeV}$.

\section{Weak interactions and flavor quantum numbers}
Hadron resonances such as $\rho$, $\omega$ or $\Delta$ have typical decay widths
$\Gamma \sim 1$--$100~\text{MeV}$, characteristic of strong decays.
However, the lightest isospin multiplets $(\pi^\pm,\pi^0,p,n)$ have much smaller widths:
\begin{itemize}
  \item $\pi^0$ decays electromagnetically with $\Gamma \sim 10^{-5}~\text{MeV}$,
  \item $\pi^\pm$ and $n$ decay weakly with $\Gamma$ as small as $10^{-14}$--$10^{-24}~\text{MeV}$.
\end{itemize}
Weak interactions are also needed to describe muon decay and strange-hadron decays.

\subsection{Strangeness}
Strange hadrons are produced in strong interactions as particle--antiparticle pairs
but decay via the weak interaction. This motivates an approximately conserved quantum number
in strong processes: \textbf{strangeness} $S$.
Strange mesons and baryons populate isospin multiplets, and the lightest ones decay weakly.

\section{Discrete symmetry violation and CP}
Parity (P) and charge conjugation (C) are maximally violated in weak interactions.
A further, smaller violation of CP is observed in neutral kaon decays. This will play a central
role in flavor physics and in understanding the origin of matter--antimatter asymmetry.

\section{More generations: charm, bottom and top}
The existence of the charm quark was predicted (around 1970) because flavor-changing neutral
currents were not observed. Charm was discovered in 1974 through the resonance
$e^+e^-\to \gamma^\ast \to J/\psi$, interpreted as a $c\bar c$ bound state.

Bottom was discovered in 1977 through $e^+e^-\to \gamma^\ast\to \Upsilon$, interpreted as a
$b\bar b$ bound state. The electroweak theory then implied a third-generation partner, the top quark,
and the bottom--top structure is crucial for understanding CP violation in kaon (and later $B$) systems.

The top quark was discovered in 1995 at the Tevatron ($p\bar p$ collider), with $m_t\simeq 172~\text{GeV}$
and a large width $\Gamma_t\sim 1.4~\text{GeV}$. It decays weakly before hadronizing, making it unique among quarks.

\section{Summary and ``big picture''}
Particle physics is organized around:
\begin{itemize}
  \item the energy-scale dependence of effective degrees of freedom,
  \item QFT as the consistent framework combining QM and SR,
  \item gauge symmetries and their associated forces,
  \item confinement and color singlets in QCD,
  \item weak interactions, flavor, and discrete symmetry violation.
\end{itemize}
These themes will guide the construction of concrete field theories and the analysis of real processes
in the next chapters.